\begin{enumerate}
	\item On sait que : 

\begin{description}[itemsep=1em]
	\item $\dfrac{8}{4}=2$ ;
	\item $\dfrac{3}{4}=0,75$ ;
	\item $\dfrac{4}{3}\approx 1,33333$ ;
	\item $1+\dfrac{3}{100}=1,03$.
\end{description}

On a donc : $0,7 < \dfrac{3}{4} < 1 < 1+\dfrac{3}{100} < 1,33 < \dfrac{4}{3} < 1,5 < 2 $.
	
	\item \begin{enumerate}
	\item Cette expérience a 2 issues favorables : $\dfrac{8}{4}=2$ et $1$. 
	
	\medskip
	Sachant que le nombre total d'issues est 8 et que chaque issue est équiprobable, la probabilité que le nombre choisi soit un entier est donc : $\dfrac{2}{8}=\dfrac{1}{4}=0,25$.

	\item Seul $\dfrac{4}{3}$ n'est pas un nombre décimal donc la probabilité que le nombre choisi soit un nombre décimal est : $\dfrac{7}{8}=0,875$.
	
	\end{enumerate}

\end{enumerate}